\section*{Capítulo \ref{ch:b3:origin-of-elements} --- De dónde proceden los elementos}
% Values from the Book 5 scratch calculator calc/ch30.py.

\begin{solution}{exo:b3:origin-of-elements:1}
(a)~$\mathrm{Si/Fe} = 10^{7.51 - 7.46} = 10^{0.05} = 1.12$ y
$\mathrm{O/Fe} = 10^{8.69 - 7.46} = 10^{1.23} = 17.0$: aproximadamente un átomo de
silicio y diecisiete átomos de oxígeno por cada átomo de hierro. (b)~$\mathrm{Fe/Au} =
10^{7.46 - 0.91} = 10^{6.55} = \num{3.55e6}$ y $\mathrm{Pb/Au} =
10^{1.95 - 0.91} = 10^{1.04} = 11.0$: más allá del pico del hierro las abundancias
caen en un factor de millones, pero el plomo, en su propio pico, es once veces más
abundante que el oro.
\end{solution}

\begin{solution}{exo:b3:origin-of-elements:2}
(a)~$10^{7.46 - 12} = \num{2.884e-5}$, es decir, \num{28.8} átomos de hierro por
millón de átomos de hidrógeno. (b)~Cien veces menos: \num{0.288} por
millón, un átomo de hierro por cada $1/\num{2.884e-7} = \num{3.47e6}$ átomos de
hidrógeno. (c)~$[\mathrm{Fe/H}] = \log_{10}(\num{e-7}/\num{2.884e-5}) =
-2.46$.
\end{solution}

\begin{solution}{exo:b3:origin-of-elements:3}
$^{124}\mathrm{Xe}$ y $^{126}\mathrm{Xe}$ se sitúan en el lado rico en protones,
fuera de la vía s, y quedan apantallados por el $^{124}\mathrm{Te}$ y
$^{126}\mathrm{Te}$, ambos estables: \omterm{def:b3:origin-of-elements:p-process}{núcleos p}. $^{128}\mathrm{Xe}$ y $^{130}\mathrm{Xe}$
están en la vía s, pero quedan protegidos del \omterm{def:b3:origin-of-elements:s-process}{proceso r} por los
$^{128}\mathrm{Te}$ y $^{130}\mathrm{Te}$, estables: solo s. $^{129}\mathrm{Xe}$,
$^{131}\mathrm{Xe}$ y $^{132}\mathrm{Xe}$ están en la vía s y no tienen
isóbaro estable de menor $Z$ (el yodo solo tiene el 127): ambos procesos.
$^{134}\mathrm{Xe}$ y $^{136}\mathrm{Xe}$ se sitúan más allá de
$^{133}\mathrm{Xe}$, donde la vía s deja el xenón, y no quedan
apantallados: solo r.
\end{solution}

\begin{solution}{exo:b3:origin-of-elements:4}
(a)~$t = d/c = 40\times\qty{3.086e22}{m}/\qty{2.998e8}{m.s^{-1}} =
\qty{4.12e15}{s}$, aproximadamente \num{130} millones de años. (b)~$\Delta v/c \le
1.7/\num{4.12e15} = \num{4.1e-16}$. Es una cota superior, ya que parte o
todo el retraso puede deberse a la emisión de los rayos gamma tras la
fusión.
\end{solution}

\begin{solution}{exo:b3:origin-of-elements:5}
(a)~$295\times0.1067 = \qty{31.48}{mb}$, $832\times0.0373 =
\qty{31.03}{mb}$, $155\times0.1984 = \qty{30.75}{mb}$. (b)~Desde
$^{122}\mathrm{Te}$: $N_{123} = 31.48/832 = \num{0.03783}$, \qty{1.4}{\percent}
por encima del valor medido, y $N_{124} = 31.48/155 = \num{0.2031}$,
\qty{2.4}{\percent} por encima. Las abundancias que difieren en un factor de cinco se
predicen con un margen de unos pocos por ciento solo a partir de las secciones eficaces.
(c)~$^{128}\mathrm{Te}$ es solo r: el flujo s deja el teluro en
$^{127}\mathrm{Te}$, que se desintegra, de modo que ningún flujo s estacionario atraviesa
$^{128}\mathrm{Te}$ y su abundancia quedó fijada por el \omterm{def:b3:origin-of-elements:s-process}{proceso r}.
\end{solution}

\begin{solution}{exo:b3:origin-of-elements:6}
(a)~$\tau_0\langle\sigma\rangle = 0.306\times155 = 47.4$, luego
$f = 1/(1 + 1/47.4) = \num{0.979}$ para $^{124}\mathrm{Te}$; para
$^{138}\mathrm{Ba}$, $0.306\times4.0 = 1.224$ y $f = 1/(1 + 0.817) =
\num{0.550}$: el núcleo mágico transmite solo el \qty{55}{\percent} del
flujo. (b)~Los seis factores son \num{0.959}, \num{0.550}, \num{0.908},
\num{0.771}, \num{0.972} y \num{0.915}; su producto es \num{0.328}:
$\sigma N$ cae en un factor de \num{3.05}. (c)~La caída medida es
$27.6/7.72 = 3.57$. El modelo acierta el tamaño del salto con un margen
del \qty{15}{\percent}, casi todo ello procedente de $^{138}\mathrm{Ba}$ y
$^{140}\mathrm{Ce}$, los dos núcleos con $N = 82$.
\end{solution}

\begin{solution}{exo:b3:origin-of-elements:7}
(a)~$kT = 8\times10^3\times\qty{1.602e-19}{J} = \qty{1.282e-15}{J}$, y
$v_T = \sqrt{2\times\num{1.282e-15}/\num{1.6749e-27}} =
\qty{1.24e6}{m.s^{-1}}$. (b)~Con $\langle\sigma\rangle =
\qty{155e-31}{m^2}$, $\lambda = n_\mathrm{n}\langle\sigma\rangle v_T$
es \qty{1.92e-11}{s^{-1}} para $n_\mathrm{n} = \qty{e12}{m^{-3}}$
($10^6$ por centímetro cúbico): un tiempo medio de \qty{5.2e10}{s}, aproximadamente
\num{1650} años; después \num{16.5} años a $10^8$ y \num{0.165}~año,
dos meses, a \qty{e10}{cm^{-3}}. (c)~$^{127}\mathrm{Te}$ se desintegra en
horas, mucho más rápido que cualquiera de estos tiempos de captura: siempre se desintegra y
la vía no se bifurca allí. $^{134}\mathrm{Cs}$ vive aproximadamente dos
años, entre los tiempos de captura a $10^8$ y \qty{e10}{cm^{-3}}: a
las altas densidades de los destellos de helio, una buena parte de él se captura
antes de desintegrarse. Es un punto de ramificación, y las abundancias a uno y otro lado
de él miden la densidad de neutrones.
\end{solution}

\begin{solution}{exo:b3:origin-of-elements:8}
(a)~$10^{6.30 - 3.14} = 1445$ veces menos abundante que el calcio y
$10^{4.97 - 3.14} = 68$ veces menos abundante que el titanio. (b)~$10^{8.46 - 1.38}
= \num{1.2e7}$. (c)~El escandio muestra el efecto par-impar: los núcleos con
número par de protones están más firmemente ligados y se forman por la
captura de núcleos de helio en las etapas de quemado de las estrellas masivas, lo que
los favorece. El berilio muestra el valle de abundancia del litio, el berilio y el boro:
estos núcleos frágiles se destruyen a las temperaturas del interior estelar,
y el quemado de helio pasa directamente del helio al carbono
sin producirlos.
\end{solution}

\begin{solution}{exo:b3:origin-of-elements:9}
(a)~Sea $\alpha$ el conjunto de los \omterm{def:b3:origin-of-elements:alpha}{elementos alfa} y $\mathrm{Fe}_\mathrm{cc}$,
$\mathrm{Fe}_\mathrm{Ia}$ el hierro procedente de los dos tipos de supernova. En
la meseta, $\alpha/\mathrm{Fe}_\mathrm{cc} = 10^{P}$ veces la
razón solar. Para $[\mathrm{Fe/H}] = 0$, $\alpha/(\mathrm{Fe}_\mathrm{cc} +
\mathrm{Fe}_\mathrm{Ia})$ es solar. Dividiendo,
$\mathrm{Fe}_\mathrm{cc}/(\mathrm{Fe}_\mathrm{cc} + \mathrm{Fe}_\mathrm{Ia})
= 10^{-P}$, y la fracción de tipo~Ia es $1 - 10^{-P} = 1 - 10^{-0.4} =
\num{0.602}$. (b)~$10^{-P} = 1/3$ da $P = \log_{10}3 = 0.48$.
(c)~Las supernovas de tipo~Ia comienzan tras el mismo retardo en ambas galaxias.
Cuando comenzaron en la galaxia enana, su gas solo había alcanzado
$[\mathrm{Fe/H}] = -2$, frente a $-1$ en la Galaxia: había convertido su gas
en estrellas diez veces más lentamente.
\end{solution}

\begin{solution}{exo:b3:origin-of-elements:10}
(a)~Para el circonio,
\[
\frac{(X/\mathrm{Mg})_\mathrm{Earth}}{(X/\mathrm{Mg})_\mathrm{CI}} =
\frac{10.5/228000}{3.96/95870} = 1.11 ;
\]
Del mismo modo $(2670/228000)/(5010/95870) = 0.224$ para el sodio,
\num{0.190} para el potasio y \num{0.118} para el rubidio. (b)~El circonio, el más refractario, se sitúa
cerca de su proporción condrítica. Los tres metales alcalinos moderadamente volátiles,
que condensan a \qtyrange{800}{1006}{K}, están empobrecidos en factores de
cinco a nueve, el rubidio, el más volátil, en mayor medida. La
Tierra silicatada sigue la \omterm{def:b3:origin-of-elements:volatility}{tendencia de volatilidad}. (c)~El magnesio es un elemento mayor,
refractario que condensa junto con la mayor parte de la roca, y es
litófilo: permaneció en el manto cuando se formó el núcleo. Un referente siderófilo
como el hierro o el níquel mezclaría la pérdida hacia el núcleo con la
pérdida por volatilidad.
\end{solution}

\begin{solution}{exo:b3:origin-of-elements:11}
(a)~Cuando una masa $\mathrm{d}s$ queda atrapada en estrellas, el viento elimina
$\lambda\,\mathrm{d}s$ de gas de \omterm{def:b3:origin-of-elements:metallicity}{metalicidad} $Z$:
$\mathrm{d}M_\mathrm{g} = -(1 + \lambda)\,\mathrm{d}s$ y
$\mathrm{d}(ZM_\mathrm{g}) = y\,\mathrm{d}s - Z\,\mathrm{d}s -
\lambda Z\,\mathrm{d}s$. Desarrollando la segunda y restando
$Z\,\mathrm{d}M_\mathrm{g}$ deja $M_\mathrm{g}\,\mathrm{d}Z =
y\,\mathrm{d}s = -\frac{y}{1 + \lambda}\,\mathrm{d}M_\mathrm{g}$, lo que
integrado da el resultado. (b)~$\ln(1/0.29) = 1.238$: $Z = 0.0155$ para
$\lambda = 0$, y la mitad, $0.0077$, para $\lambda = 1$. (c)~Un viento
reduce el rendimiento efectivo en el factor $1 + \lambda$. En una galaxia de
masa pequeña, las supernovas expulsan el gas con más facilidad, $\lambda$ es grande,
y los metales se pierden con el gas.
\end{solution}

\begin{solution}{exo:b3:origin-of-elements:12}
(a)~$N_\mathrm{Li}/N_\mathrm{H} = 10^{0.96 - 12} = \num{9.1e-12}$ en el
Sol y $10^{3.25 - 12} = \num{1.78e-9}$ en las condritas CI. El Sol
conserva $10^{-2.29} = \num{0.0051}$ de su litio: ha destruido
\qty{99.5}{\percent} de él. (b)~$\num{1.78e-9}/\num{1.45e-10} = 12.3$ y
$\num{1.78e-9}/\num{4.72e-10} = 3.8$. (c)~El Sistema Solar contiene casi
cuatro veces más litio que el que produjo el Big Bang: la mayor parte de su litio se
formó más tarde, por algún otro proceso. Las estrellas más antiguas retienen tres veces
menos de lo que produjo el Big Bang: o bien han destruido parte de él, como
ha hecho el Sol, o bien la predicción omite algo; las dos comparaciones
por sí solas no dicen cuál.
\end{solution}

\begin{solution}{pb:b3:origin-of-elements:1}

\textbf{Parte I.}

\textbf{1.} $10^{12 - 7.46} = 10^{4.54} = \num{34700}$.

\textbf{2.} $N_\mathrm{He}/N_\mathrm{H} = 10^{10.914 - 12} = \num{0.0820}$, así que
$4N_\mathrm{He}/N_\mathrm{H} = \num{0.328}$, frente a $Y/X =
0.2423/0.7438 = \num{0.326}$. Coinciden con una diferencia de \qty{0.7}{\percent}; la
pequeña diferencia se debe a que un átomo de helio pesa algo menos que cuatro
átomos de hidrógeno.

\textbf{3.} $\log_{10}\bigl((1/350000)/\num{2.884e-5}\bigr) = -5.544 + 4.540 =
-1.00$.

\textbf{4.} $[\mathrm{O/H}] = -1.00 + 0.40 = -0.60$, así que
$\log\varepsilon(\mathrm{O}) = 8.69 - 0.60 = 8.09$. La estrella tiene
$10^{8.09 - 6.46} = 42.7$ átomos de oxígeno por átomo de hierro, frente a
$10^{1.23} = 17.0$ en el Sol: $10^{0.4} = 2.5$ veces más.

\textbf{5.} $Z \approx 0.0139\times10^{-0.60} = \num{0.0035}$, una cuarta parte de la
del Sol.

\textbf{6.} El hierro tiene muchísimas líneas en el espectro de una estrella fría, mientras que
el oxígeno tiene pocas y débiles: el hierro se mide mucho más fácilmente y
con mayor precisión, y representa al conjunto de los metales.

\textbf{Parte II.}

\textbf{7.} Cada uno se sitúa tras un isóbaro estable de menor $Z$:
$^{122}\mathrm{Sn}$, $^{123}\mathrm{Sb}$ y $^{124}\mathrm{Sn}$ detienen
las cadenas de desintegración del \omterm{def:b3:origin-of-elements:s-process}{proceso r} antes de que alcancen el telurio. Los tres se
encuentran en la vía s, que recorre el telurio desde $^{122}\mathrm{Te}$ hasta
$^{126}\mathrm{Te}$.

\textbf{8.} \num{31.48}, \num{31.03} y \qty{30.75}{mb}; media
\qty{31.09}{mb}.

\textbf{9.} La media de las otras dos es $(31.48 + 30.75)/2 = \qty{31.11}{mb}$,
así que $N_{123} = 31.11/832 = \num{0.03740}$, solo \qty{0.3}{\percent}
por encima del valor medido \num{0.0373}.

\textbf{10.} $1.329\times44.4 = \qty{59.0}{mb}$, casi el doble de la meseta.
$^{128}\mathrm{Te}$ es solo r: la vía s abandona el telurio en
$^{127}\mathrm{Te}$, así que su abundancia se formó por el \omterm{def:b3:origin-of-elements:s-process}{proceso r} y
no debe nada al flujo s.

\textbf{11.} $1 - 7.43/8.53 = \num{0.13}$: \qty{13}{\percent} del flujo pasó
eludió $^{148}\mathrm{Sm}$, mediante capturas sobre los núcleos
inestables $^{147}\mathrm{Nd}$, $^{147}\mathrm{Pm}$ y $^{148}\mathrm{Pm}$ antes de
que pudieran desintegrarse en samario.

\textbf{12.} Con una capa cerrada de 82 neutrones, estos núcleos ligan
débilmente un neutrón adicional y tienen secciones eficaces de captura muy pequeñas (\qty{4.0}{mb}
para $^{138}\mathrm{Ba}$). El flujo se acumula sobre ellos y solo una parte
lo atraviesa: son cuellos de botella.

\textbf{Parte III.}

\textbf{13.} $\mu_1 = 13.7/(13.7 + 33.4) = 13.7/47.1 = \num{0.291}$.

\textbf{14.} $\ln(1/0.291) = \num{1.235}$, así que $y = 0.0154/1.235 = \num{0.0125}$,
\num{0.81} veces $Z_\odot$.

\textbf{15.} Con $\mu = 0.9$, $\ln(1/0.9) = 0.1054$, así que $Z = 0.0125\times0.1054 =
\num{0.00131}$. Como $0.00131/0.0154 = 0.0851$, $[\mathrm{Fe/H}] =
-1.07$. Con $\mu = 0.5$, $Z = 0.0125\times0.693 = \num{0.00864}$ y
$[\mathrm{Fe/H}] = -0.25$.

\textbf{16.} $Z = 0.00154$, así que $\mu = \mathrm{e}^{-0.00154/0.01247} =
\mathrm{e}^{-0.1235} = \num{0.884}$.

\textbf{17.} $1 - 0.884 = \num{0.116}$: \qty{12}{\percent} de la masa.

\textbf{18.} $y = 0.0154/\ln 10 = \num{0.0067}$, aproximadamente la mitad.

\textbf{Parte IV.}

\textbf{19.} La masa de las estrellas con \omterm{def:b3:origin-of-elements:metallicity}{metalicidad} inferior a $Z$ es
$M(1 - \mathrm{e}^{-Z/y})$, de modo que por unidad de $Z$ vale
$(M/y)\,\mathrm{e}^{-Z/y}$. Como $Z = Z_\odot10^{[\mathrm{Fe/H}]}$,
$\mathrm{d}Z = Z\ln 10\,\mathrm{d}[\mathrm{Fe/H}]$, y por unidad de
$[\mathrm{Fe/H}]$ la distribución es proporcional a
$Z\mathrm{e}^{-Z/y}$. Su derivada,
$(1 - Z/y)\mathrm{e}^{-Z/y}$, se anula en $Z = y$: el máximo está en
$[\mathrm{Fe/H}] = \log_{10}0.810 = -0.09$.

\textbf{20.} $\num{0.116}/(1 - 0.291) = 0.1162/0.7091 = \num{0.164}$.

\textbf{21.} $Z/y = 1.235\times10^{-0.5} = 0.3905$, $1 - \mathrm{e}^{-0.3905} =
0.3233$, y $0.3233/0.7091 = \num{0.456}$.

\textbf{22.} $6670\times0.1638 = 1093$, frente a 26: unas cuarenta veces más.

\textbf{23.} $26/902 = \qty{2.9}{\percent}$, todavía más de cinco veces por debajo
de la caja. Las estrellas del modelo son las de vida larga, las de como mucho
una masa solar aproximadamente, que han sobrevivido todas desde que se formaron;
entre las estrellas F más masivas, las más antiguas, que son también las más
pobres en metales, ya han muerto, y la muestra completa las subcuenta.

\textbf{24.} $(1 - 0.1^{0.1})/(1 - 0.1) = 0.2057/0.9 = \num{0.23}$: peor.
Aumentar la fracción de gas tampoco ayuda: a medida que $\mu_1$ tiende a 1,
la fracción tiende a $0.1$, todavía más de tres veces el
\qty{2.9}{\percent} observado. Ninguna fracción de gas salva la caja cerrada.

\textbf{25.} La caja cerrada predice que
\textbf{\qty[text-series-to-math=true]{16}{\percent}} de las estrellas de vida larga
de la vecindad solar tienen $[\mathrm{Fe/H}] < -1$, frente a un \qty{3}{\percent}
como máximo observado. El disco no fue una caja cerrada: creció por una lenta
caída de gas pobre en metales, de modo que casi nunca hubo gas esperando a baja
\omterm{def:b3:origin-of-elements:metallicity}{metalicidad} para formar estrellas.
\end{solution}
